\begin{helper}[GnpGraphWeightNaturalDegreeConcentration]
Fix \(\varepsilon>0\), let \(\beta=1/2\), and let
\(m:\mathbb N\to\mathbb N\) satisfy
\[
m(n)=\noderef{TwoBiteNaturalReducedVertexCount}(n)
\]
for every \(n\).  Put
\[
p(n)=\noderef{TwoBiteEdgeProbability}(\beta,n).
\]
Under the \(G(m(n),p(n))\) product graph law encoded by
\(\noderef{GnpGraphWeight}\), the probability that some vertex has degree
outside relative error \(\varepsilon\) from \(p(n)m(n)\) is negligible:
\[
\sum_G
\begin{cases}
\noderef{GnpGraphWeight}(p(n),G),
&\text{if there is }v\in\operatorname{Fin}(m(n))\text{ such that}\\
&\quad\neg\noderef{WithinRelativeError}
  \!\left(\varepsilon,\ p(n)m(n),\
    \noderef{GraphDegreeCount}(G,v)\right),\\
0,&\text{otherwise}
\end{cases}
\longrightarrow 0 .
\]
\end{helper}
\begin{proof}
Write \(M=m(n)\), \(p=p(n)\), and \(L=\log n\).  We first record the
asymptotic input.  By \(\noderef{TwoBiteNaturalReducedVertexCount}\),
together with its defining support \(\noderef{TwoBiteReducedVertexCount}\)
and \(\noderef{TwoBiteStretch}\),
\[
M=\left\lfloor \frac{n}{L^2}\right\rfloor .
\]
For all sufficiently large \(n\), \(L\ge1\) and
\((n/L^2)\ge2\).  Hence the elementary floor estimate
\(\lfloor x\rfloor\ge x/2\), valid for \(x\ge2\), gives
\[
\frac12\frac{n}{L^2}\le M\le \frac{n}{L^2}.
\tag{1}
\]
Since \(\beta=1/2\), \(\noderef{TwoBiteEdgeProbability}\) gives
\[
p=\frac12\sqrt{\frac{L}{n}}.
\tag{2}
\]
For all large \(n\), also \(0\le p\le1\): nonnegativity follows from
\(L\ge0\), and \(p\le1\) follows from \(\log n\le n\).  Equations (1) and
(2) imply
\[
pM\ge \frac14\frac{\sqrt n}{L^{3/2}}\to\infty,
\qquad
\log M=O(L),
\qquad
\frac{pM}{\log M}\to\infty .
\tag{3}
\]
Indeed \(M\le n\), so \(\log M\le L\) once \(M\ge1\), and powers of \(n\)
dominate powers of \(\log n\).  Discarding finitely many initial \(n\)'s
does not affect the \(\noderef{NegligibleProbability}\) conclusion.

We next derive the one-vertex degree distribution from the finite product
law.  Let
\[
E_M=\{(a,b)\in \operatorname{Fin}(M)\times\operatorname{Fin}(M):a<b\}
\]
be the ordered representative set of unordered graph edges used in
\(\noderef{GnpGraphWeight}\).  The definition
\(\noderef{GnpGraphWeight}\) says that, for a graph \(G\),
\[
\noderef{GnpGraphWeight}(p,G)
=\prod_{e=(a,b)\in E_M}
  \begin{cases}
  p,&G.Adj(a,b),\\
  1-p,&\neg G.Adj(a,b).
  \end{cases}
\tag{4}
\]
Each factor is nonnegative because \(0\le p\le1\), and
\(\noderef{GnpGraphWeightSumOne}\) gives total mass \(1\) when all edge
coordinates are summed.

For this finite coordinate set, summing over graphs is the same as summing
over all Boolean assignments to \(E_M\): a graph gives the assignment
recording which represented edges are adjacent, and an assignment gives the
unique simple graph whose adjacency relation is the symmetric loopless
relation generated by the present represented edges.  This is the coordinate
bijection used in \(\noderef{GnpGraphWeightSumOne}\).

For large \(n\) we have \(M\ge2\).  Fix \(v\in\operatorname{Fin}(M)\).
Split the edge-coordinate set into
\[
I_v=\{e\in E_M:e\text{ is incident to }v\},\qquad
N_v=E_M\setminus I_v .
\]
Since the graph is simple, there is no loop at \(v\), and each
\(u\ne v\) contributes exactly one coordinate \(\{u,v\}\).  Thus
\(|I_v|=M-1\).  By \(\noderef{GraphDegreeCount}\),
\[
D_v=\noderef{GraphDegreeCount}(G,v)
\]
is exactly the number of incident coordinates in \(I_v\) that are present.

Let \(t\in\mathbb R\).  We compute the moment generating function of
\(D_v\) directly from (4):
\[
\sum_G e^{tD_v(G)}\noderef{GnpGraphWeight}(p,G).
\]
Summing first over the nonincident coordinates \(N_v\), each nonincident
coordinate contributes
\[
p+(1-p)=1.
\]
The finite product version of this collapse is precisely
\(\noderef{ProductWeightSumOne}\); it is the same product-sum calculation
used globally in \(\noderef{GnpGraphWeightSumOne}\), now applied only to the
coordinates in \(N_v\).  Therefore all nonincident coordinates contribute
total factor \(1\).  For an incident coordinate \(e\in I_v\), the two
possibilities contribute
\[
(1-p)\cdot e^{t\cdot0}+p\cdot e^{t\cdot1}=1-p+pe^t.
\]
Multiplying over the \(M-1\) incident coordinates gives the exact identity
\[
\sum_G e^{tD_v(G)}\noderef{GnpGraphWeight}(p,G)
=(1-p+pe^t)^{M-1}.
\tag{5}
\]
This is the full finite-product justification for treating \(D_v\) as a
binomial sum; no independence is being assumed beyond the product
factorization in \(\noderef{GnpGraphWeight}\).

The mean obtained from (5) is
\[
\mu=(M-1)p=pM-p.
\tag{6}
\]
Equivalently, it is the sum over the \(M-1\) incident coordinates of their
individual present-edge mass \(p\).

Let \(\delta=\min\{\varepsilon/4,1/2\}\).  For a binomial sum \(Y\) with mean
\(\mu\), the standard Chernoff moment calculation gives constants
\(c_+(\delta)>0\) and \(c_-(\delta)>0\) such that
\[
\Pr(Y\ge (1+\delta)\mu)\le e^{-c_+(\delta)\mu},
\qquad
\Pr(Y\le (1-\delta)\mu)\le e^{-c_-(\delta)\mu}.
\tag{7}
\]
Here is the calculation using (5).  For \(\lambda>0\),
\[
\Pr(D_v\ge(1+\delta)\mu)
\le
\exp(-\lambda(1+\delta)\mu)
(1-p+pe^\lambda)^{M-1}.
\]
Since \(1+x\le e^x\),
\[
(1-p+pe^\lambda)^{M-1}
\le\exp(\mu(e^\lambda-1)).
\]
Taking \(\lambda=\log(1+\delta)\) gives exponent
\[
-\mu\bigl((1+\delta)\log(1+\delta)-\delta\bigr),
\]
and \((1+\delta)\log(1+\delta)-\delta>0\).  This proves the upper-tail
bound with
\[
c_+(\delta)=(1+\delta)\log(1+\delta)-\delta.
\]
For the lower tail, set \(\lambda=-\log(1-\delta)>0\).  On the event
\(D_v\le(1-\delta)\mu\), one has
\(e^{-\lambda D_v}\ge e^{-\lambda(1-\delta)\mu}\), so Markov's inequality
and (5) with \(t=-\lambda\) give
\[
\Pr(D_v\le(1-\delta)\mu)
\le
\exp(\lambda(1-\delta)\mu)
(1-p+pe^{-\lambda})^{M-1}.
\]
Again using \(1+x\le e^x\) and \(e^{-\lambda}=1-\delta\), the exponent is
\[
-\mu\bigl(\delta+(1-\delta)\log(1-\delta)\bigr).
\]
The quantity \(\delta+(1-\delta)\log(1-\delta)\) is positive for
\(0<\delta<1\), so this proves the lower-tail bound with
\[
c_-(\delta)=\delta+(1-\delta)\log(1-\delta)>0.
\]
Let \(c=\min\{c_+(\delta),c_-(\delta)\}>0\).  Combining the two tails,
\[
\sum_G
\begin{cases}
\noderef{GnpGraphWeight}(p,G),
& |D_v(G)-\mu|>\delta\mu,\\
0,&\text{otherwise}
\end{cases}
\le 2e^{-c\mu}.
\tag{8}
\]

For all sufficiently large \(n\), \(p\le(\varepsilon/4)pM\) and
\(\mu\ge pM/2\), by (3) and (6).  If
\(|D_v-\mu|\le\delta\mu\), then
\[
|D_v-pM|
\le |D_v-\mu|+|\mu-pM|
\le \delta\mu+p
\le \frac{\varepsilon}{4}pM+\frac{\varepsilon}{4}pM
\le \varepsilon pM.
\]
Thus failure of
\(\noderef{WithinRelativeError}(\varepsilon,pM,D_v)\) is contained in the
two-sided Chernoff tail in (8), and hence, for all large \(n\),
\[
\sum_G
\begin{cases}
\noderef{GnpGraphWeight}(p,G),
&\neg\noderef{WithinRelativeError}(\varepsilon,pM,D_v(G)),\\
0,&\text{otherwise}
\end{cases}
\le 2e^{-c pM/2}.
\tag{9}
\]

There are \(M\) possible vertices.  The graph weights are nonnegative, so
the finite union bound gives
\[
\sum_G
\begin{cases}
\noderef{GnpGraphWeight}(p,G),
&\exists v\in\operatorname{Fin}(M):
  \neg\noderef{WithinRelativeError}(\varepsilon,pM,D_v(G)),\\
0,&\text{otherwise}
\end{cases}
\le 2M e^{-c pM/2}.
\tag{10}
\]
Taking logarithms of the right-hand side gives
\[
\log 2+\log M-\frac c2 pM\to-\infty
\]
by (3).  Therefore the right-hand side of (10) tends to \(0\), which is
exactly the displayed \(\noderef{NegligibleProbability}\) statement for the
bad degree event under the \(\noderef{GnpGraphWeight}\) product law.
\end{proof}
